%%%%%%%%%%%%%%%%%%%%%%%%%%%%%%%%%%%%%%%%%%%%%%%%%% 
% Picture with arrows - mathrm
%%%%%%%%%%%%%%%%%%%%%%%%%%%%%%%%%%%%%%%%%%%%%%%%%% 

\usetikzlibrary{arrows}

% Define how TiKZ will draw the nodes
\tikzset{mathterm/.style={draw=white,fill=white,rectangle,anchor=base}}
\tikzstyle{every picture}+=[remember picture]
\everymath{\displaystyle}

\makeatletter

%%%%%%%%%%%%%%%%%%%%%%%%%%%%%%%%%%%%%%%%%%%%%%%%%% 
% Designate a term in a math environment to point to
% Syntax: \mathterm[node label]{some math}
%%%%%%%%%%%%%%%%%%%%%%%%%%%%%%%%%%%%%%%%%%%%%%%%%% 

\newcommand\mathterm[2][]{%
   \tikz [baseline] { \node [mathterm] (#1) {$#2$}; }}

%%%%%%%%%%%%%%%%%%%%%%%%%%%%%%%%%%%%%%%%%%%%%%%%%% 
% A command to draw an arrow from the current 
% position to a labelled math term
% Default color=black, default arrow head=stealth
% Syntax: \indicate[color]{term to point to}[path options]
%%%%%%%%%%%%%%%%%%%%%%%%%%%%%%%%%%%%%%%%%%%%%%%%%% 

\newcommand\indicate[2][black]{%
   \tikz [baseline] \node [inner sep=0pt,anchor=base] (i#2) {\vphantom|};
   \@ifnextchar[{\@indicateopts{#1}{#2}}{\@indicatenoopts{#1}{#2}}}
\def\@indicatenoopts#1#2{%
   {\color{#1} \tikz[overlay] \path[line width=1pt,draw=#1,-stealth] (i#2) edge (#2);}}
\def\@indicateopts#1#2[#3]{%
   {\color{#1} \tikz[overlay] \path[line width=1pt,draw=#1,-stealth] (i#2) [#3] edge (#2);}}

\makeatother

%%%%%%%%%%%%%%%%%%%%%%%%%%%%%%%%%%%%%%%%%%%%%%%%%% 
% Footnote for title
%%%%%%%%%%%%%%%%%%%%%%%%%%%%%%%%%%%%%%%%%%%%%%%%%% 

\let\ACMmaketitle=\maketitle
\renewcommand{\maketitle}{\begingroup\let\footnote=\thanks \ACMmaketitle\endgroup}

\def\tang{\ThisStyle{\abovebaseline[0pt]{\scalebox{-1}{$\SavedStyle\perp$}}}}

%%%%%%%%%%%%%%%%%%%%%%%%%%%%%%%%%%%%%%%%%%%%%%%%%% 
%
% Insert coding
%
% Example: 
% \lstinputlisting[language=Octave, numbers=none]{code/fsine.m}
%%%%%%%%%%%%%%%%%%%%%%%%%%%%%%%%%%%%%%%%%%%%%%%%%% 

\definecolor{codegreen}{rgb}{0,0.6,0}
\definecolor{codegray}{rgb}{0.5,0.5,0.5}
\definecolor{codepurple}{rgb}{0.58,0,0.82}
\definecolor{backcolor}{rgb}{0.95,0.95,0.92}
\definecolor{mygreen}{RGB}{28,172,0}

\lstdefinestyle{mystyle}{
    %backgroundcolor=\color{backcolor},   
    commentstyle=\color{mygreen},
    keywordstyle=\color{blue},
    numberstyle= \tiny\color{codegray},
    stringstyle=\color{codepurple},
    basicstyle=\ttfamily\footnotesize,
    %frame=L,
    %frame = lines,
    %frame = single,
    breakatwhitespace=false,         
    breaklines=true,                 
    captionpos=b,                    
    keepspaces=true,
    %numbers=left, 
    %numbersep=7pt,                  
    showspaces=false,                
    showstringspaces=false,
    showtabs=false,                  
    tabsize=2,
    % xleftmargin=5pt,
    % framexleftmargin=10pt,
    % framextopmargin=4pt,
    % framexbottommargin=4pt,
    %emph=[1]{for,end,break},emphstyle=[1]\color{blue},
    %emph=[2]{length,gca}, emphstyle=[2]\color{black}
}
 
\lstset{style=mystyle}

%%%%%%%%%%%%%%%%%%%%%%%%%%%%%%%%%%%%%%%%%%%%%%%%%%%%%
\newcounter{mycode}
\colorlet{colexam}{green!75!black}
\tcbset{
  base/.style={
    breakable,
    enhanced,
    frame engine = path,
    sharp corners,
    title,
    attach boxed title to top left,
    boxed title style={
            size=minimal, top=0pt, left=0pt
            },
    coltitle=black,
    fonttitle=\bfseries,
  }
}
%
\newtcolorbox[use counter=mycode]{code}{%
  base,
  boxed title style={
  overlay={
      \node[anchor=north east,inner ysep=0pt,align=right,text width=2in] 
        at ([yshift=-0.55ex]frame.north west) {\hfill {In [\thetcbcounter]:}};
  }},
  borderline west={3pt}{0pt}{colexam},
  borderline north={1pt}{0pt}{gray!30},
  borderline south={1pt}{0pt}{gray!30},
  borderline east={1pt}{0pt}{gray!30},
  frame hidden,
  colback = gray!5,
}

%%%%%%%%%%%%%%%%%%%%%%%%%%%%%%%%%%%%%%%%%%%%%%%%%% 
% Volume in Fluid
%%%%%%%%%%%%%%%%%%%%%%%%%%%%%%%%%%%%%%%%%%%%%%%%%% 

\makeatletter
\DeclareRobustCommand{\volume}{\text{\volumedash}V}
\newcommand{\volumedash}{%
  \makebox[0pt][l]{%
    \ooalign{\hfil\hphantom{$\m@th V$}\hfil\cr\kern0.08em\textbf{--}\hfil\cr}%
  }%
}

%%%%%%%%%%%%%%%%%%%%%%%%%%%%%%%%%%%%%%%%%%%%%%%%%% 
% CREATING NICE RED BOX
% https://tex.stackexchange.com/questions/719491/how-can-i-get-the-title-in-tcolorbox-partly-framed-and-not-framed
%%%%%%%%%%%%%%%%%%%%%%%%%%%%%%%%%%%%%%%%%%%%%%%%%% 

% Defining counter:
\newcounter{defcounter}
\setcounter{defcounter}{0}
\renewcommand{\thedefcounter}{\arabic{defcounter}}

\newtcolorbox{definition}[1][]{%
    enhanced,
    before skip=5mm, after skip=5mm, 
    colback=white, colframe=red, % Colors
    rounded corners, arc=3mm,
    drop fuzzy shadow, % Shadow
    fonttitle=\bfseries,
    title={#1}
}


%%%%%%%%%%%%%%%%%%%%%%%%%%%%%%%%%%%%%%%%%%%%%%%%%%
% Circuits
%%%%%%%%%%%%%%%%%%%%%%%%%%%%%%%%%%%%%%%%%%%%%%%%%%
\usepackage[european,s traightvoltages, RPvoltages, americanresistor, americaninductors]{circuitikz}
\tikzset{every picture/.style={line width=.2mm}}

% Tikz Library
\usetikzlibrary{calc}

% Bipoles Specifications
\ctikzset{bipoles/thickness=2, label distance=1mm, voltage shift = 1}

% Arrows Above Compenents
% Source: https://tex.stackexchange.com/questions/574576/circuitikz-straight-voltage-arrows-with-fixed-length
\newcommand{\fixedvlen}[3][0.4cm]{% [semilength]{node}{label}
    % get the center of the standard arrow
    \coordinate (#2-Vcenter) at ($(#2-Vfrom)!0.5!(#2-Vto)$);
    % draw an arrow of a fixed size around that center and on the same line
    \draw[-{Triangle[round,open]}] ($(#2-Vcenter)!#1!(#2-Vfrom)$) -- ($(#2-Vcenter)!#1!(#2-Vto)$);
    % position the label as it where if standard voltages were used
    \node[ anchor=\ctikzgetanchor{#2}{Vlab}] at (#2-Vlab) {#3};
}
\newcommand{\fixedvlendashed}[3][0.75cm]{% [semilength]{node}{label}
    % get the center of the standard arrow
    \coordinate (#2-Vcenter) at ($(#2-Vfrom)!0.5!(#2-Vto)$);
    % draw an arrow of a fixed size around that center and on the same line
    \draw[dashed,-{Triangle[round,open]}] ($(#2-Vcenter)!#1!(#2-Vfrom)$) -- ($(#2-Vcenter)!#1!(#2-Vto)$);
    % position the label as it where if standard voltages were used
    \node[ anchor=\ctikzgetanchor{#2}{Vlab}] at (#2-Vlab) {#3};
}


% EXAMPLE
% \begin{circuitikz}
% %		%Grid
% %		\draw[thin, dotted] (0,0) grid (8,8);
% %		\foreach \i in {1,...,8}
% %		{
% %			\node at (\i,-2ex) {\i};	
% %		}
% %		\foreach \i in {1,...,8}
% %		{
% %			\node at (-2ex,\i) {\i};	
% %		}
% %		\node at (-2ex,-2ex) {0};
			
% 		%Coordinates
% 		\coordinate (A) at (0,8);
% 		\coordinate (B) at (0,4);
% 		\coordinate (C) at (2.5,4);
% 		\coordinate (D) at (7,4);
% 		\coordinate (E) at (2.5,0);
% 		\coordinate (Z) at (0,0);
% 		\coordinate (Earth) at (0,-1);
		
% 		%Circuit
% 		\draw (A) to[resistor, l_=$R_1$, i>=$I_{R_1}$, v^<, name=r1] (B) to[short, i>=$$] (C)
% 				 to[resistor, l^=$R_4$, a_= {$V_{R_4}= 0 = I_{R_4}$}, name=r4, *-*] (D);
% 		\draw (B) to[resistor, l_=$R_2$, i>=$I_{R_2}$, v^<, name=r2, *-*] (Z) to[short, i>=$$] (Earth) node[tlground]{};
% 		\draw (C) to[resistor, l_=$R_3$, i>=$I_{R_3}$, v^<, name=r3, *-*] (E) to[short, i>=$$] (Z);
		
% 		%Nodes
% 		\node[shift={(0.6,0)}] at (A) {$V_{CC}$};
% 		\node[left] at (A) {A};
% 		\node[left] at (B) {B};
% 		\node[above] at (C) {C};
% 		\node[above] at (D) {D};
% 		\node[right] at (E) {E};
% 		\node[left] at (Z) {F};
		
% 		% Voltages
% 		\fixedvlen[0.5cm]{r1}{$V_{R_1}$}		
% 		\fixedvlen[0.5cm]{r2}{$V_{R_2}$}	
% 		\fixedvlen[0.5cm]{r3}{$V_{R_3}$}
		
% \end{circuitikz}